% Part: computability
% Chapter: recursive-functions
% Section: introduction

\documentclass[../../../include/open-logic-section]{subfiles}

\begin{document}

\olfileid[fr]{cmp}{rec}{int}
\olsection{Introduction}

Pour élaborer une théorie mathématique de la calculabilité, il faut
d'abord en construire un \emph{modèle}. Nous voyons aujourd'hui dans
le calcul le type d'opération qu'effectuent les ordinateurs, lesquels
manipulent des symboles. Mais, aux débuts de la théorie de la
calculabilité, le calcul \emph{numérique} en était l'exemple type.
Les mathématiciens s'intéressaient depuis longtemps aux fonctions
numériques calculables, c'est-à-dire aux fonctions
$f\colon\Nat^n\to\Nat$ dont on peut calculer les valeurs. Il n'est
donc pas surprenant que la théorie de la calculabilité ait d'abord
étudié de telles fonctions. Les exemples les plus familiers,
comme l'addition, la multiplication et l'exponentiation des entiers
naturels, présentent un trait commun : on peut les définir
\emph{récursivement}. Il est donc naturel de chercher à définir la
notion générale de \emph{fonction calculable} au moyen de définitions
récursives. Parmi les nombreuses façons de définir récursivement des
fonctions numériques, un schéma particulièrement simple jouera ici
un rôle central : la \emph{récurrence primitive}.

Outre les fonctions calculables, on peut étudier les ensembles et
les relations calculables. Un ensemble est calculable si l'on peut
décider, pour tout nombre donné, s'il en est !!a{element}. Une
relation est calculable si l'on peut décider, pour tout uplet
$\tuple{n_1,\dots,n_k}$, s'il en est !!a{element}. Grâce à la
\emph{fonction caractéristique} d'un ensemble ou d'une relation,
l'étude des ensembles et des relations calculables se ramène à
celle des fonctions calculables. On peut donc également définir
les relations récursives primitives; la relation « $n$ divise
$m$ », par exemple, est récursive primitive.

Les fonctions récursives primitives, définissables à partir des fonctions de base par
composition et récurrence primitive, ne sont cependant pas les seules
fonctions numériques calculables. Parmi les diverses extensions de
ce schéma, la plus puissante et la plus communément admise consiste
à autoriser une recherche non bornée. On obtient ainsi la définition
des \emph{fonctions récursives partielles}, puis celle, apparentée,
des \emph{fonctions récursives générales}. Celles-ci sont calculables
et totales : la définition retient exactement les fonctions
récursives partielles qui se trouvent être totales. Les fonctions
récursives permettent de simuler les autres modèles de calcul,
notamment les machines de Turing et le lambda-calcul. Elles
constituent ainsi l'un des modèles de calcul reconnus.

\end{document}
